\documentclass[11pt,a4paper]{article}

\newcommand{\cabecalhodireita}{Lista de Exercícios 2}
% Preâmbulo compartilhado entre a lista de exercícios e o gabarito.
\usepackage[T1]{fontenc}
\usepackage[utf8]{inputenc}
\usepackage[portuguese]{babel}
\usepackage[a4paper,margin=2.3cm,top=2.0cm,bottom=2.2cm]{geometry}
\usepackage{amsmath,amssymb}
\usepackage{graphicx}
\usepackage{booktabs}
\usepackage{xcolor}
\usepackage{enumitem}
\usepackage{listings}
\usepackage{fancyhdr}
\usepackage[framemethod=default]{mdframed}
\usepackage{titlesec}
\usepackage{needspace}
\usepackage[hidelinks]{hyperref}

% --- Identidade visual (mesma paleta dos slides: VS Code Light+) ---
\definecolor{ufsverde}{HTML}{1F6F43}
\definecolor{bgcolor}{HTML}{F7F7F7}
\definecolor{linecolor}{HTML}{DDDDDD}
\definecolor{numcolor}{HTML}{999999}
\definecolor{kwcolor}{HTML}{0000FF}
\definecolor{strcolor}{HTML}{A31515}
\definecolor{cmtcolor}{HTML}{008000}

\lstdefinestyle{ufsStyle}{
  backgroundcolor=\color{bgcolor},
  basicstyle=\ttfamily\small\color{black},
  keywordstyle=\color{kwcolor}\bfseries,
  stringstyle=\color{strcolor},
  commentstyle=\color{cmtcolor}\itshape,
  numberstyle=\tiny\color{numcolor},
  numbers=left, numbersep=8pt,
  xleftmargin=18pt, framexleftmargin=18pt,
  breaklines=true, tabsize=4,
  keepspaces=true, showspaces=false,
  showstringspaces=false, showtabs=false,
  frame=single, rulecolor=\color{linecolor},
  framesep=4pt, captionpos=b, upquote=true,
  columns=fullflexible,
  literate={á}{{\'a}}1 {é}{{\'e}}1 {í}{{\'i}}1 {ó}{{\'o}}1 {ú}{{\'u}}1
           {â}{{\^a}}1 {ê}{{\^e}}1 {ô}{{\^o}}1 {ã}{{\~a}}1 {õ}{{\~o}}1
           {ç}{{\c c}}1 {Á}{{\'A}}1 {É}{{\'E}}1 {Í}{{\'I}}1 {Ó}{{\'O}}1
           {Ú}{{\'U}}1 {Â}{{\^A}}1 {Ê}{{\^E}}1 {Ô}{{\^O}}1 {Ã}{{\~A}}1
           {Õ}{{\~O}}1 {Ç}{{\c C}}1,
}
\lstset{style=ufsStyle, language=C}

% Atalho para código no meio do texto
\newcommand{\cd}[1]{\texttt{#1}}

% --- Seções ---
\titleformat{\section}{\normalfont\large\bfseries\color{ufsverde}}{}{0pt}{}
\titlespacing*{\section}{0pt}{14pt}{6pt}

% --- Cabeçalho e rodapé ---
\pagestyle{fancy}
\fancyhf{}
\renewcommand{\headrulewidth}{0.4pt}
\renewcommand{\footrulewidth}{0pt}
\fancyhead[L]{\footnotesize COMP0512 --- Programação C $\cdot$ T02 $\cdot$ 2026.2}
\fancyhead[R]{\footnotesize\textbf{\cabecalhodireita}}
\fancyfoot[C]{\footnotesize\thepage}

% --- Ambiente de exercício ---
\newcounter{exnum}
\newcommand{\exercicio}[1]{%
  \par\Needspace{8\baselineskip}\refstepcounter{exnum}%
  \vspace{9pt}%
  \noindent{\color{ufsverde}\bfseries Exercício \theexnum.}\hspace{0.4em}\textbf{#1}\par
  \vspace{3pt}%
}

% --- Bloco destacado (instruções, avisos) ---
\newmdenv[
  backgroundcolor=bgcolor,
  linecolor=linecolor,
  linewidth=0.6pt,
  leftmargin=0pt, rightmargin=0pt,
  innerleftmargin=10pt, innerrightmargin=10pt,
  innertopmargin=9pt, innerbottommargin=9pt,
  skipabove=9pt, skipbelow=9pt,
]{destaque}

% --- Cabeçalho de capa ---
\newcommand{\capa}[2]{%
  \noindent
  \begin{minipage}[c]{0.15\textwidth}
    \includegraphics[width=\textwidth]{ufs_principal_positiva.png}
  \end{minipage}%
  \hspace{0.02\textwidth}%
  \begin{minipage}[c]{0.80\textwidth}
    \raggedright\setlength{\parindent}{0pt}
    {\footnotesize Universidade Federal de Sergipe --- Departamento de Computação\par}
    {\footnotesize COMP0512 --- Programação C $\cdot$ Turma T02 $\cdot$ 2026.2\par}
    \vspace{3pt}
    {\Large\bfseries\color{ufsverde} #1\par}
    \vspace{2pt}
    {\small #2\par}
    {\footnotesize Prof.\ André Yoshiaki Kashiwabara \quad$\cdot$\quad \texttt{andreyoshiaki@dcomp.ufs.br}\par}
  \end{minipage}
  \vspace{6pt}
  \hrule height 1pt
  \vspace{10pt}
}

\setlength{\emergencystretch}{2em}

\begin{document}

\capa{Lista de Exercícios 2 --- Recursividade, Listas, Pilhas e Filas}
     {Conteúdo da segunda unidade: recursividade e pilha de execução, efeitos da
      recursão na memória, listas encadeadas simples e duplas, pilhas e filas.}

\begin{destaque}
\textbf{Entrega manuscrita.}
\begin{itemize}[leftmargin=1.2em,itemsep=1pt,topsep=3pt]
  \item Resolva \textbf{à mão}, em folhas avulsas. Escreva nome e matrícula em
        \textbf{todas} as folhas e numere-as (1/6, 2/6, \dots).
  \item Indique o número do exercício e a letra do item antes de cada resposta.
        Não é preciso copiar o enunciado.
  \item Código em C deve estar legível e indentado. Não é preciso escrever
        \cd{\#include} nem \cd{main}, a menos que o item peça.
  \item Nos desenhos de memória, use uma caixa para cada nó ou variável e uma seta
        para cada ponteiro; escreva \cd{NULL} explicitamente onde houver.
  \item Quando o item pedir ``explique'' ou ``justifique'', duas ou três frases
        bastam --- mas a resposta sem justificativa não é considerada.
\end{itemize}
\vspace{4pt}
\textbf{Data de entrega:} \rule{4cm}{0.4pt}
\end{destaque}

\noindent Salvo indicação em contrário, os exercícios usam os tipos vistos em aula:

\begin{lstlisting}
typedef struct no {            /* lista simples */
    int valor;
    struct no *prox;
} No;

typedef struct nod {           /* lista dupla */
    int valor;
    struct nod *ant, *prox;
} NoD;

typedef struct {
    NoD *ini, *fim;            /* primeiro e ultimo nos, ou NULL */
} Lista;
\end{lstlisting}

\noindent e as interfaces \cd{pilha.h} (\cd{pilha\_cria}, \cd{pilha\_libera},
\cd{empilha}, \cd{desempilha}, \cd{topo}, \cd{vazia}) e \cd{fila.h}
(\cd{fila\_cria}, \cd{fila\_libera}, \cd{enfileira}, \cd{desenfileira},
\cd{frente}, \cd{vazia}, \cd{tamanho}), todas com elementos \cd{int}.

\section{Parte I --- Recursividade e pilha de execução}

\exercicio{O que acontece na ida e o que acontece na volta.}
Considere a função abaixo.

\begin{lstlisting}
void h(int n) {
    if (n == 0) {
        printf("* ");
        return;
    }
    printf("%d ", n);
    h(n / 2);
    printf("%d ", n % 2);
}
\end{lstlisting}

\begin{enumerate}[label=(\alph*),leftmargin=1.8em,itemsep=1pt]
  \item Escreva o que \cd{h(6)} imprime e o que \cd{h(11)} imprime.
  \item Identifique o caso base e o passo recursivo. Por que a recursão sempre
        termina para $n \ge 0$?
  \item Os números impressos \textbf{depois} do \cd{*} têm um significado. Qual?
        Por que eles saem nessa ordem, e não na ordem inversa?
  \item Desenhe a pilha de execução no instante em que o \cd{*} é impresso durante
        \cd{h(6)}: um registro por chamada, com o valor de \cd{n} em cada um.
        Quantos registros de \cd{h} existem nesse instante?
  \item Quantos registros de \cd{h} ficam na pilha, no máximo, durante \cd{h(1000)}?
        Dê a conta, não só o número.
\end{enumerate}

\exercicio{Árvore de chamadas e o custo de recalcular.}
Considere a função abaixo.

\begin{lstlisting}
int t(int n) {
    if (n <= 2)
        return n;
    return t(n - 1) + t(n - 3);
}
\end{lstlisting}

\begin{enumerate}[label=(\alph*),leftmargin=1.8em,itemsep=1pt]
  \item Calcule \cd{t(3)}, \cd{t(4)}, \cd{t(5)} e \cd{t(6)}.
  \item Desenhe a árvore de chamadas de \cd{t(5)}. Quantas chamadas a \cd{t}
        acontecem, contando a primeira?
  \item Na árvore do item (b), quais chamadas são repetidas? Qual é a altura máxima
        que a pilha de execução atinge (em registros de \cd{t})?
  \item Escreva \cd{int t\_memo(int n, int *tab)} que guarda em \cd{tab[i]} cada
        valor já calculado (\cd{-1} indica ``ainda não calculado''), e uma função
        \cd{int t\_rapida(int n)} que aloca \cd{tab} no heap \textbf{uma única vez},
        inicializa, chama \cd{t\_memo}, libera e devolve o resultado.
  \item Por que a tabela deve ser alocada em \cd{t\_rapida}, e não dentro de
        \cd{t\_memo}?
\end{enumerate}

\exercicio{Escrever funções recursivas.}
Escreva as funções abaixo \textbf{sem laços}. Para cada uma, marque no código qual é
o caso base e qual é o passo recursivo.

\begin{lstlisting}
int    conta_digitos(int n);              /* n >= 0; conta_digitos(0) == 1 */
int    maior(const int *v, int n);        /* n >= 1 */
int    palindromo(const char *s, int i, int j);
double potencia(double x, int n);         /* n >= 0 */
\end{lstlisting}

\begin{enumerate}[label=(\alph*),leftmargin=1.8em,itemsep=1pt]
  \item Implemente \cd{conta\_digitos} e \cd{maior}.
  \item Implemente \cd{palindromo}, que devolve 1 se \cd{s[i..j]} se lê igual nos
        dois sentidos e 0 caso contrário. Mostre a chamada inicial para uma string
        \cd{s} inteira. Quais são os \textbf{dois} casos base?
  \item Implemente \cd{potencia} usando $x^n = (x^{n/2})^2$ para $n$ par e
        $x^n = x \cdot x^{n-1}$ para $n$ ímpar, de modo que cada chamada faça
        \textbf{uma só} chamada recursiva. Quantas chamadas acontecem para
        $n = 16$? E se você escrevesse \cd{potencia(x, n/2) * potencia(x, n/2)}?
  \item Um colega escreveu a função abaixo. O que acontece com \cd{soma(6)}? E com
        \cd{soma(5)}? Corrija o caso base.
\end{enumerate}

\begin{lstlisting}
int soma(int n) {          /* deveria somar n + (n-2) + (n-4) + ... */
    if (n == 0)
        return 0;
    return n + soma(n - 2);
}
\end{lstlisting}

\section{Parte II --- Recursividade e memória}

\exercicio{Locais que se multiplicam pela profundidade.}
Considere a função abaixo. Suponha uma pilha de execução de 8\,MiB
($8 \times 2^{20}$ bytes) e que cada registro de ativação de \cd{niveis} ocupe
4\,KiB (4096 bytes), quase todos do vetor \cd{linha}.

\begin{lstlisting}
void niveis(int n) {
    char linha[4000];
    if (n == 0)
        return;
    snprintf(linha, sizeof linha, "nivel %d", n);
    puts(linha);
    niveis(n - 1);
}
\end{lstlisting}

\begin{enumerate}[label=(\alph*),leftmargin=1.8em,itemsep=1pt]
  \item Aproximadamente, qual é o maior \cd{n} para o qual \cd{niveis(n)} termina
        sem estourar a pilha? Mostre a conta.
  \item Se \cd{linha} fosse declarado com \cd{malloc(4000)} dentro da função (e
        liberado antes do \cd{return}), esse limite mudaria? Explique o que passa a
        ficar na pilha e o que passa a ficar no heap.
  \item Reescreva a função para que o vetor seja alocado \textbf{uma única vez} e
        passado para as chamadas seguintes como parâmetro. Quem aloca e quem libera?
  \item A chamada \cd{niveis(n - 1)} é uma chamada de cauda? Um compilador com
        \cd{-O2} poderia evitar o estouro. Por que, ainda assim, você não deve
        contar com isso?
\end{enumerate}

\exercicio{Recursão de cauda.}
Considere as quatro funções abaixo.

\begin{lstlisting}
int A(int n)          { if (n == 0) return 0;   return 1 + A(n - 1); }
int B(int n, int acc) { if (n == 0) return acc; return B(n - 1, acc + n); }
void C(int n)         { if (n == 0) return; printf("%d ", n); C(n - 1); }
int D(int n)          { if (n == 0) return 1;   return 2 * D(n - 1); }
\end{lstlisting}

\begin{enumerate}[label=(\alph*),leftmargin=1.8em,itemsep=1pt]
  \item Diga o que \cd{A(4)}, \cd{B(4, 0)} e \cd{D(4)} devolvem e o que \cd{C(4)}
        imprime.
  \item Classifique cada uma como recursão de cauda ou não. Para as que não são,
        diga exatamente \textbf{o que fica pendente} em cada registro de ativação.
  \item Reescreva \cd{A} e \cd{D} em recursão de cauda, usando um parâmetro
        acumulador. Mostre a chamada inicial de cada uma.
  \item Reescreva \cd{B} com um laço \cd{while}, sem recursão. O que os parâmetros
        \cd{n} e \cd{acc} viraram?
\end{enumerate}

\exercicio{Quem recebe o bloco é dono dele.}
A função abaixo devolve, num bloco do heap, a string com o caractere \cd{c} repetido
\cd{k} vezes.

\begin{lstlisting}
char *repete(char c, int k) {
    if (k == 0) {
        char *s = malloc(1);
        s[0] = '\0';
        return s;
    }
    char *resto = repete(c, k - 1);
    char *s = malloc(k + 1);
    s[0] = c;
    strcpy(s + 1, resto);
    return s;
}
\end{lstlisting}

\begin{enumerate}[label=(\alph*),leftmargin=1.8em,itemsep=1pt]
  \item Para \cd{repete('x', 4)}, quantas chamadas a \cd{malloc} acontecem e quantos
        bytes são alocados no total?
  \item O chamador faz \cd{char *s = repete('x', 4); puts(s); free(s);}.
        Quantos blocos e quantos bytes ficam perdidos? Generalize para \cd{k}.
  \item Corrija a função com o menor número de linhas possível. Em que ponto exato
        o \cd{free} entra, e por que não pode vir antes do \cd{strcpy}?
  \item Na versão corrigida, quantos bytes de heap ficam ocupados \textbf{ao mesmo
        tempo}, no máximo, durante \cd{repete('x', k)}? Compare com a versão
        original.
  \item A função não confere o retorno de \cd{malloc}. Acrescente a verificação de
        modo que, em caso de falha, ela devolva \cd{NULL} sem vazar nada.
\end{enumerate}

\section{Parte III --- Listas encadeadas simples}

\exercicio{Seguir as setas.}
Considere o trecho abaixo, em que \cd{cria(v)} devolve um nó novo com valor \cd{v} e
\cd{prox} igual a \cd{NULL}.

\begin{lstlisting}
No *a = cria(1), *b = cria(2), *c = cria(3), *d = cria(4);
a->prox = c;  c->prox = b;  b->prox = d;
No *ini = a;
                                   /* bloco 1 */
No *p = ini->prox;
ini->prox = p->prox;
p->prox = ini;
ini = p;
                                   /* bloco 2 */
No *q = ini;
while (q->prox->prox != NULL)
    q = q->prox;
q->prox->prox = ini;
ini = q->prox;
q->prox = NULL;
\end{lstlisting}

\begin{enumerate}[label=(\alph*),leftmargin=1.8em,itemsep=1pt]
  \item Desenhe a lista (a partir de \cd{ini}) antes do bloco 1, depois do bloco 1 e
        depois do bloco 2. Escreva, para cada desenho, os valores na ordem em que
        um percurso a partir de \cd{ini} os visitaria.
  \item Descreva em uma frase o que o bloco 1 faz e o que o bloco 2 faz, para uma
        lista qualquer.
  \item Ao final do laço do bloco 2, quem é \cd{q}? Por que a condição testa
        \cd{q->prox->prox}, e não \cd{q->prox}?
  \item O que acontece com o bloco 2 se a lista tiver um único nó? E se estiver
        vazia? Acrescente as proteções necessárias.
\end{enumerate}

\exercicio{Funções sobre uma lista simples.}
Escreva as funções abaixo. Nenhuma delas deve imprimir nada.

\begin{lstlisting}
int  conta(const No *ini, int x);        /* quantas vezes x aparece        */
int  ordenada(const No *ini);            /* 1 se nao decrescente, 0 se nao */
No  *remove_todos(No *ini, int x);       /* tira todos os x; novo inicio   */
void tira_repetidos(No *ini);            /* lista ordenada: deixa um de cada */
\end{lstlisting}

\begin{enumerate}[label=(\alph*),leftmargin=1.8em,itemsep=1pt]
  \item Implemente \cd{conta} e \cd{ordenada}. O que \cd{ordenada} devolve para a
        lista vazia e para a lista de um nó? Justifique a escolha.
  \item Implemente \cd{remove\_todos}, liberando cada nó removido. Mostre que ela
        funciona para a lista \cd{7 -> 7 -> 3 -> 7 -> NULL} com \cd{x = 7}:
        desenhe a lista depois de cada nó removido.
  \item Implemente \cd{tira\_repetidos}: \cd{1 -> 1 -> 2 -> 3 -> 3 -> 3 -> NULL}
        vira \cd{1 -> 2 -> 3 -> NULL}. Por que ela não precisa devolver um novo
        início, ao contrário de \cd{remove\_todos}?
  \item Diga o custo de cada uma das quatro funções em função do número $n$ de nós.
\end{enumerate}

\exercicio{O \cd{insere\_fim} que não insere.}
Um colega escreveu a função abaixo para inserir \cd{x} no fim da lista.

\begin{lstlisting}
void insere_fim(No *ini, int x) {
    No *novo = malloc(sizeof(No));
    novo->valor = x;
    No *p = ini;
    while (p != NULL)
        p = p->prox;
    p->prox = novo;
}
\end{lstlisting}

\begin{enumerate}[label=(\alph*),leftmargin=1.8em,itemsep=1pt]
  \item A função tem \textbf{quatro} defeitos. Liste-os, dizendo em que linha está
        cada um e o que ele causa em tempo de execução.
  \item Mesmo depois de corrigir o laço, a chamada \cd{insere\_fim(lista, 5)} com
        \cd{lista == NULL} não altera \cd{lista} na \cd{main}. Por quê?
  \item Reescreva a função de duas formas corretas: (1)
        \cd{No *insere\_fim(No *ini, int x)}, devolvendo o novo início, e (2)
        \cd{void insere\_fim2(No **pini, int x)}. Mostre a chamada de cada uma.
  \item Quanto custa inserir $n$ valores, um a um, no fim de uma lista que começa
        vazia, usando essa função? O que mudaria na estrutura para que cada inserção
        no fim custasse um tempo constante?
\end{enumerate}

\exercicio{Ponteiro para a seta.}
Considere a função abaixo e a lista \cd{10 -> 20 -> 30 -> NULL}, guardada em
\cd{No *ini} na \cd{main}.

\begin{lstlisting}
void remove_valor(No **pp, int x) {
    while (*pp != NULL && (*pp)->valor != x)
        pp = &(*pp)->prox;
    if (*pp != NULL) {
        No *morto = *pp;
        *pp = morto->prox;
        free(morto);
    }
}
\end{lstlisting}

\begin{enumerate}[label=(\alph*),leftmargin=1.8em,itemsep=1pt]
  \item Mostre a chamada correta a partir da \cd{main}.
  \item Para \cd{x = 10}, \cd{x = 30} e \cd{x = 99}, diga para \textbf{que variável
        ou campo} \cd{pp} aponta quando o laço termina (por exemplo: ``a variável
        \cd{ini} da \cd{main}'' ou ``o campo \cd{prox} do nó 20'').
  \item Desenhe a memória no caso \cd{x = 30}, logo antes do \cd{free}: os três nós,
        \cd{ini}, \cd{pp} e \cd{morto}.
  \item Por que essa versão não precisa de um \cd{if} separado para o primeiro nó,
        ao contrário da versão com \cd{ant}?
  \item O que aconteceria se as duas condições do \cd{while} fossem escritas na
        ordem inversa? Para qual dos três valores de \cd{x} o problema aparece?
\end{enumerate}

\section{Parte IV --- Listas duplamente encadeadas}

\exercicio{Quatro setas para mover um nó.}
A lista \cd{l} contém \cd{10 <-> 20 <-> 30 <-> 40}, com \cd{l->ini} no 10 e
\cd{l->fim} no 40. O ponteiro \cd{p} aponta para o nó 20 e \cd{q} para o nó 40.

\begin{lstlisting}
p->ant->prox = p->prox;     /* 1 */
p->prox->ant = p->ant;      /* 2 */
p->ant  = q;                /* 3 */
p->prox = NULL;             /* 4 */
q->prox = p;                /* 5 */
l->fim  = p;                /* 6 */
\end{lstlisting}

\begin{enumerate}[label=(\alph*),leftmargin=1.8em,itemsep=1pt]
  \item Desenhe a lista depois da linha 2 e depois da linha 6, com todas as setas
        \cd{ant} e \cd{prox}, \cd{l->ini} e \cd{l->fim}.
  \item Escreva o que um percurso de \cd{l->ini} pelos \cd{prox} imprime e o que um
        percurso de \cd{l->fim} pelos \cd{ant} imprime, ao final.
  \item Diga o valor de \cd{l->ini->prox->prox->ant->valor} ao final.
  \item Se \cd{p} apontasse para o nó 10 (o primeiro), que linha falharia, e por quê?
        Reescreva as linhas 1 e 2 para funcionarem também nesse caso.
  \item E se \cd{p} e \cd{q} apontassem ambos para o 40? Siga as linhas na ordem e diga
        o que acontece.
\end{enumerate}

\exercicio{Operações nos extremos e no meio.}
Escreva as funções abaixo para a lista dupla com \cd{ini} e \cd{fim}.

\begin{lstlisting}
void insere_depois(Lista *l, NoD *p, int x);  /* p pertence a l */
int  remove_ultimo(Lista *l);                 /* devolve o valor; lista nao vazia */
void imprime_reversa(const Lista *l);
\end{lstlisting}

\begin{enumerate}[label=(\alph*),leftmargin=1.8em,itemsep=1pt]
  \item Implemente \cd{insere\_depois}. Quantos ponteiros mudam quando \cd{p} está
        no meio? E quando \cd{p == l->fim}?
  \item Implemente \cd{remove\_ultimo} sem laço. Não esqueça o caso em que a lista
        tem um único nó: o que deve acontecer com \cd{ini} e \cd{fim}?
  \item Implemente \cd{imprime\_reversa} sem recursão e sem estrutura auxiliar.
  \item Remover o último nó de uma lista \textbf{simples} com \cd{ini} e \cd{fim}
        custa quanto? Por que manter o ponteiro \cd{fim} não resolve o problema na
        lista simples, mas resolve na dupla?
\end{enumerate}

\exercicio{O nó sentinela.}
Uma lista dupla circular com nó cabeça \cd{cab} (que não guarda valor) tem
\cd{cab->prox} no primeiro nó e \cd{cab->ant} no último; a lista vazia tem
\cd{cab->ant == cab->prox == cab}. A inserção é:

\begin{lstlisting}
void insere_antes(NoD *p, int x) {
    NoD *novo = malloc(sizeof *novo);
    novo->valor = x;
    novo->prox = p;
    novo->ant  = p->ant;
    p->ant->prox = novo;
    p->ant = novo;
}
\end{lstlisting}

\begin{enumerate}[label=(\alph*),leftmargin=1.8em,itemsep=1pt]
  \item Desenhe a lista vazia. Depois, desenhe a lista após a chamada
        \cd{insere\_antes(cab, 5)} e, em seguida, após \cd{insere\_antes(cab, 7)}.
  \item As duas chamadas do item (a) inseriram no início ou no fim? Que chamada
        insere no outro extremo?
  \item Escreva \cd{int tamanho(const NoD *cab)}. Qual é a condição de parada do
        percurso, e por que não pode ser \cd{p != NULL}?
  \item Escreva \cd{void remove\_no(NoD *p)} em três linhas, sem nenhum \cd{if}.
        Explique por que a versão sem sentinela precisava de dois \cd{if}.
  \item O que acontece se alguém chamar \cd{remove\_no(cab)}? Como a função
        poderia se proteger, se recebesse também \cd{cab}?
\end{enumerate}

\section{Parte V --- Pilhas}

\exercicio{Seguir a pilha.}
Responda usando apenas a definição de pilha (o último a entrar é o primeiro a sair).

\begin{enumerate}[label=(\alph*),leftmargin=1.8em,itemsep=1pt]
  \item Avalie a expressão pós-fixa \cd{5 1 2 + 4 * + 3 -} com uma pilha: número
        empilha; operador desempilha dois valores e empilha o resultado. Desenhe a
        pilha depois de cada símbolo e dê o resultado final.
  \item No item (a), quando o operador \cd{-} desempilha dois valores, qual deles é
        o operando da esquerda? Mostre o que sairia se a ordem fosse trocada.
  \item Os valores \cd{1 2 3 4} são empilhados nessa ordem, mas a qualquer momento
        pode-se desempilhar e imprimir o topo. Para cada sequência abaixo, diga se
        ela pode ser impressa; se puder, dê a sequência de \cd{empilha} e
        \cd{desempilha}; se não puder, diga em que momento ela trava.\\[2pt]
        \cd{2 4 3 1} \qquad \cd{3 1 4 2} \qquad \cd{4 3 2 1} \qquad \cd{1 4 2 3}
\end{enumerate}

\exercicio{Pilha num vetor que dobra.}
Um colega implementou a pilha num vetor com a estrutura abaixo.

\begin{lstlisting}
struct pilha {
    int *v;
    int n;      /* quantos elementos ha; o topo e v[n-1] */
    int cap;    /* quantos cabem em v */
};

void empilha(Pilha *p, int x) {
    if (p->n == p->cap) {
        p->cap = 2 * p->cap;
        p->v = realloc(p->v, p->cap * sizeof *p->v);
    }
    p->v[p->n++] = x;
}

int desempilha(Pilha *p) {
    return p->v[p->n--];
}

int topo(const Pilha *p) {
    return p->v[p->n];
}
\end{lstlisting}

\begin{enumerate}[label=(\alph*),leftmargin=1.8em,itemsep=1pt]
  \item \cd{desempilha} e \cd{topo} têm o mesmo erro. Qual é? Com a pilha contendo
        \cd{10 20 30} (30 no topo), o que cada uma devolveria?
  \item Nenhuma das duas trata a pilha vazia. O que acontece se \cd{desempilha} for
        chamada com \cd{n == 0}, já com o erro do item (a) corrigido?
  \item \cd{empilha} tem dois problemas no \cd{realloc}: um quando ele falha e outro
        quando a capacidade inicial é 0. Explique ambos e reescreva a função.
  \item Começando com \cd{cap = 1}, empilham-se 1000 valores. Quantas vezes o vetor
        cresce? Quantos elementos são copiados no total pelos \cd{realloc} (no pior
        caso, em que cada um muda o bloco de lugar)? E se, em vez de dobrar, a
        capacidade aumentasse de 1 em 1?
\end{enumerate}

\exercicio{Pilha numa lista e o uso só pela interface.}
\begin{enumerate}[label=(\alph*),leftmargin=1.8em,itemsep=1pt]
  \item Escreva \cd{struct pilha}, \cd{empilha}, \cd{desempilha} e
        \cd{pilha\_libera} para uma pilha guardada numa lista simples (use o tipo
        \cd{No}).
  \item O topo deve ficar no \textbf{primeiro} ou no \textbf{último} nó? Justifique
        com o custo de \cd{empilha} e \cd{desempilha} em cada escolha.
  \item Usando \textbf{apenas} as funções de \cd{pilha.h} (sem acessar os campos),
        escreva \cd{int base(Pilha *p)}, que devolve o elemento do fundo de uma
        pilha não vazia e deixa a pilha exatamente como estava. De que estrutura
        auxiliar você precisou?
  \item A função do item (c) funciona sem nenhuma alteração tanto com a pilha num
        vetor quanto com a pilha numa lista. Que característica de \cd{pilha.h}
        garante isso?
\end{enumerate}

\section{Parte VI --- Filas}

\exercicio{A fila circular dá a volta.}
Uma fila circular guarda \cd{v}, \cd{ini} (posição da frente), \cd{n} (quantos há) e
\cd{cap}. Ela começa vazia, com \cd{cap = 5} e \cd{ini = 0}; o valor novo entra em
\cd{v[(ini + n) \% cap]}.

\begin{lstlisting}
enfileira(f, 10);  enfileira(f, 20);  enfileira(f, 30);  enfileira(f, 40);
desenfileira(f);   desenfileira(f);   desenfileira(f);
enfileira(f, 50);  enfileira(f, 60);  enfileira(f, 70);
\end{lstlisting}

\begin{enumerate}[label=(\alph*),leftmargin=1.8em,itemsep=1pt]
  \item Desenhe \cd{v[0..4]} e dê \cd{ini} e \cd{n} depois da primeira linha, depois
        da segunda e depois da terceira. Marque as posições que guardam valores que
        já não pertencem à fila.
  \item Em que posição entra o próximo valor? Escreva a conta.
  \item Agora \cd{enfileira(f, 80)} e \cd{enfileira(f, 90)}. O 90 não cabe e a fila
        cresce para \cd{cap = 10}, copiando os elementos na ordem da fila para o
        início do vetor novo. Desenhe \cd{v[0..9]}, \cd{ini} e \cd{n} depois do 90.
  \item Por que não basta fazer \cd{v = realloc(v, 10 * sizeof *v)}, como na pilha?
        Em que ordem os valores sairiam nesse caso?
  \item Uma versão alternativa guarda \cd{ini} e \cd{fim} (próxima posição livre) em
        vez de \cd{ini} e \cd{n}. Depois do \cd{enfileira(f, 80)}, quanto valem
        \cd{ini} e \cd{fim}? Por que, com só esses dois campos, não dá para
        distinguir a fila cheia da fila vazia?
\end{enumerate}

\exercicio{Fila numa lista: as duas pontas.}
Um colega implementou a fila numa lista simples com \cd{struct fila \{ No *ini, *fim; int n; \};}.

\begin{lstlisting}
void enfileira(Fila *f, int x) {
    No *novo = malloc(sizeof *novo);
    novo->valor = x;
    novo->prox = NULL;
    f->fim->prox = novo;
    f->fim = novo;
    f->n++;
}

int desenfileira(Fila *f) {
    No *velho = f->ini;
    int x = velho->valor;
    f->ini = velho->prox;
    free(velho);
    f->n--;
    return x;
}
\end{lstlisting}

\begin{enumerate}[label=(\alph*),leftmargin=1.8em,itemsep=1pt]
  \item O que acontece na primeira chamada de \cd{enfileira} sobre uma fila recém
        criada, com \cd{ini} e \cd{fim} iguais a \cd{NULL}? Corrija a função
        testando se \cd{f->fim} é \cd{NULL}.
  \item Com \cd{enfileira} corrigida, faça estas três chamadas, nesta ordem:\\
        \cd{enfileira(f, 1);} \quad \cd{desenfileira(f);} \quad \cd{enfileira(f, 2);}\\
        Desenhe \cd{ini}, \cd{fim} e os nós depois de cada uma. Onde está o defeito de \cd{desenfileira} e qual
        categoria de erro de memória ele produz?
  \item Corrija \cd{desenfileira}, incluindo o tratamento da fila vazia.
  \item Por que a fila sai pelo \textbf{primeiro} nó e entra depois do
        \textbf{último}, e não o contrário? Diga o custo de cada operação nas duas
        escolhas.
\end{enumerate}

\exercicio{Resolver problemas com filas e pilhas.}
\begin{enumerate}[label=(\alph*),leftmargin=1.8em,itemsep=1pt]
  \item Na brincadeira da batata quente, $n$ crianças numeradas de 1 a $n$ estão
        numa fila; a cada rodada, a da frente passa para o fim $k$ vezes e, em
        seguida, a que ficou na frente sai. Para $n = 6$ e $k = 2$, desenhe a fila
        no início de cada rodada e dê a ordem em que as crianças saem.
  \item Usando apenas \cd{fila.h}, escreva
        \cd{void intercala(Fila *a, Fila *b, Fila *saida)}, que tira alternadamente
        um elemento de \cd{a} e um de \cd{b} e os coloca em \cd{saida}; quando uma
        acaba, o resto da outra vai inteiro.
  \item Usando uma fila e uma pilha (um \cd{char} cabe num \cd{int}), escreva
        \cd{int eh\_palindromo(const char *s)}. Explique por que as duas estruturas,
        juntas, resolvem o problema.
  \item \cd{pilha.h} e \cd{fila.h}, como estão, não podem ser incluídos no mesmo
        arquivo \cd{.c}. Por quê? Proponha uma correção nas duas interfaces.
\end{enumerate}

\exercicio{Escolher a estrutura.}
\begin{enumerate}[label=(\alph*),leftmargin=1.8em,itemsep=1pt]
  \item Copie e preencha a tabela com o custo de cada operação ($O(1)$ ou $O(n)$)
        numa coleção de $n$ elementos.
\end{enumerate}

\begin{center}
\small
\begin{tabular}{lcccc}
\toprule
Operação & Vetor & Lista simples (\cd{ini}) & Lista simples (\cd{ini}, \cd{fim}) & Lista dupla (\cd{ini}, \cd{fim}) \\
\midrule
inserir no início   & & & & \\
inserir no fim      & & & & \\
remover do início   & & & & \\
remover do fim      & & & & \\
acessar o $k$-ésimo & & & & \\
buscar um valor     & & & & \\
\bottomrule
\end{tabular}
\end{center}

\begin{enumerate}[label=(\alph*),leftmargin=1.8em,itemsep=1pt,start=2]
  \item Na coluna ``Vetor'', considere que ``inserir no fim'' usa um vetor que dobra
        quando enche. Por que se diz que essa operação custa $O(1)$
        \emph{amortizado}, se de vez em quando ela copia o vetor inteiro?
  \item Para cada situação abaixo, diga se você usaria uma pilha, uma fila, uma lista
        simples ou uma lista dupla, e justifique em uma frase:
        \begin{enumerate}[label=(\roman*),leftmargin=1.6em,itemsep=0pt]
          \item o ``desfazer'' de um editor de texto;
          \item os documentos esperando numa impressora compartilhada;
          \item uma \emph{playlist} com os botões ``anterior'' e ``próxima'' e
                remoção da música atual;
          \item verificar se as \cd{\{ \}} de um programa estão balanceadas;
          \item as senhas de atendimento de uma agência bancária.
        \end{enumerate}
  \item Para qual das situações do item (c) mais de uma estrutura serve bem?
        Compare as duas opções.
\end{enumerate}

\end{document}
